%% kapa: exam=1 round=19 year=2008 period=1 form=A subject=civil_law subject_ko=민법 seq=23 number=23
%% source: Q-Net 감정평가사 기출문제 (cst00309) artlSeq=1001962 file=extracted/1008/1001962/2008년 제19회 감정평가사 1차시험 기출문제/1교시.hwp
%% conversion: kordoc 4.24.51 -> tools/gen1.py (auto); review=proofread-sonnet; flags=none
\begin{question}{23}
다음 설명 중 옳은 것은? (다툼이 있는 경우에는 판례에 의함)
\begin{choices}
  \item 현실의 인도가 되려면 양도인의 물건에 대한 사실상의 지배가 동일성을 유지한 채 양수인에게 전면적․배타적으로 이전되어야 한다.
  \item 건물의 증축부분이 구조상․이용상의 독립성을 갖추면 구분소유를 설정하려는 소유자의 의사에 관계없이 구분소유권이 성립한다.
  \item 무권리자로부터 점유개정에 의하여 동산의 점유를 취득한 경우에도 선의취득을 인정할 수 있다.
  \item 부동산의 소유자가 부합에 의하여 동산의 소유권을 취득한 경우에는 부동산소유자는 동산소유자에 대하여 보상할 의무가 없다.
  \item 무효․취소를 가지고 선의의 제3자에게 대항할 수 없는 경우에도 예고등기가 허용된다.
\end{choices}
\end{question}
\answer{1}
